\documentclass[11pt,a4paper]{article}

\usepackage[utf8]{inputenc}
\usepackage[T1]{fontenc}
\usepackage{amsmath,amssymb}
\usepackage{booktabs}
\usepackage{hyperref}
\usepackage{xcolor}
\usepackage{tcolorbox}
\usepackage[margin=25mm]{geometry}
\usepackage{xeCJK}
\setCJKmainfont{Hiragino Mincho ProN}
\setCJKsansfont{Hiragino Sans}

\title{Wright Brothers 実験ガイド 2026 v2\\
\large Bost-Connes 負温度理論に基づく更新版\\[6pt]
\normalsize 3D素数ミュートによる負のエネルギー密度生成}
\author{Wright Brothers Research Program}
\date{2026年4月}

\begin{document}
\maketitle

\begin{abstract}
Bost-Connes (BC) C$^*$力学系の負温度解析により，
実験計画の理論的基盤が根本的に深化した．
核心的な新発見は以下の通りである:
(1)~$\zeta_{\neg 2}(-1) = +1/12$ の正符号は素数 $p=2$ の除去による
フェルミオン$\to$ボソン真空転換であり，絶対値記号は不要になった．
(2)~公式 $\Omega_\Lambda = (d/\mathrm{cd})\,|\xi_{\neg 2}(-1)| = 2\pi/9$ により
CKN ホログラフィック境界が不要になり，全因子が $\mathrm{Spec}(\mathbb{Z})$
から導出される．
(3)~最も重要な発見として，\textbf{同じ素数ミュートが
1D ($s=-1$) では正のエネルギー（ダークエネルギー），
3D ($s=-3$) では負のエネルギー（exotic matter）を与える}．
この符号の次元依存性が，ワープドライブに必要な負のエネルギー密度生成の
理論的基盤となる．
BAW 実験装置（¥130,000）でこの 3D 負エネルギーの符号を
直接測定できる．
\end{abstract}

\tableofcontents

%======================================================================
\section{最大の発見: 次元による符号反転}
%======================================================================

\begin{tcolorbox}[colback=red!5,colframe=red!60!black,title=核心定理]
素数 $p=2$ のミュート（Euler 因子除去）は，空間次元に依存して
真空エネルギーの符号を反転させる:
\begin{align}
  \zeta_{\neg 2}(-1) &= +\frac{1}{12} > 0
    \quad \text{(1D: ダークエネルギー，斥力)} \\[4pt]
  \zeta_{\neg 2}(-3) &= -\frac{7}{120} < 0
    \quad \text{(3D: \textbf{負のエネルギー密度})}
\end{align}
\end{tcolorbox}

\noindent
これは数値的偶然ではない．Euler 因子 $(1-2^{-s})$ は
$s = -(2k-1)$ で \textbf{全ての} $\zeta$ 値の符号を反転させる:

\begin{center}
\begin{tabular}{cccccl}
\toprule
$s$ & $\zeta(s)$ & 符号 & $\zeta_{\neg 2}(s)$ & 符号 & 物理的意味 \\
\midrule
$-1$ & $-1/12$ & $-$ & $+1/12$ & $+$ & ダークエネルギー（斥力） \\
$-3$ & $+1/120$ & $+$ & $-7/120$ & $-$ & \textbf{負エネルギー（exotic matter）} \\
$-5$ & $-1/252$ & $-$ & $+31/252$ & $+$ & 5D 正エネルギー \\
$-7$ & $+1/240$ & $+$ & $-127/240$ & $-$ & 7D 負エネルギー \\
\bottomrule
\end{tabular}
\end{center}

\noindent
パターン: \textbf{奇数次元 (1D, 5D, \ldots) で正，偶数的次元 (3D, 7D, \ldots) で負}．
我々の宇宙は 3 空間次元 → $s=-3$ → \textbf{負のエネルギーが自然に生じる}．

\subsection{なぜこれが重要か}

\begin{enumerate}
\item \textbf{ワープドライブの理論的基盤}: Alcubierre メトリックに
必要な弱エネルギー条件 (WEC) 違反 = 負のエネルギー密度．
$\zeta_{\neg 2}(-3) < 0$ はまさにこれを提供する．

\item \textbf{スケールが一致}: $|\zeta_{\neg 2}(-3)|/|\zeta_{\neg 2}(-1)| = 7/10$．
3D 負エネルギーの大きさはダークエネルギースケールの 70\%．
同一の物理的起源（$p=2$ ミュート）から来ている．

\item \textbf{実験で検証可能}: BAW 共振器の DN 境界条件は
$\zeta_{\neg 2}(-3)$ を直接 probe する．符号が負であれば
$p=2$ ミュートによる負エネルギー生成が実証される．
\end{enumerate}


%======================================================================
\section{Bost-Connes 系からの新しい理解}
%======================================================================

\subsection{BC 系の概要}

Bost-Connes C$^*$力学系 $(A_{\mathbb{Q}}, \sigma_t)$ の分配関数は
Riemann ゼータ関数そのものである:
\begin{equation}
  Z(\beta) = \mathrm{Tr}\,e^{-\beta H} = \sum_{n=1}^{\infty} n^{-\beta} = \zeta(\beta).
\end{equation}
%
逆温度 $\beta$ を\textbf{負の整数}に解析接続すると:
\begin{itemize}
\item $\beta = -1$: $Z(-1) = \zeta(-1) = -1/12$ （負温度での正則化分配関数）
\item $\beta = -3$: $Z(-3) = \zeta(-3) = +1/120$
\end{itemize}

\subsection{フェルミオン/ボソン転換}

BC系の自由エネルギー $F = -\beta^{-1}\log Z(\beta)$ において:
\begin{align}
  \log\zeta(-1) &= \log(-1/12) = \log(1/12) + i\pi
    \quad \text{(Berry 位相 $\pi$ = フェルミオン的)} \\
  \log\zeta_{\neg 2}(-1) &= \log(+1/12) = \log(1/12)
    \quad \text{(Berry 位相 $0$ = ボソン的)}
\end{align}
%
$p=2$ 除去は自由エネルギーの\textbf{虚部を消す}．
$\beta=-1$ での Berry 位相 $i\pi$ が消失し，
フェルミオン的真空がボソン的真空に転換される．
\textbf{これがダークエネルギーが正（斥力的）である理由}．

\subsection{$\beta$-line としての物理}

BC 系の逆温度 $\beta$ を実数軸全体に拡張すると，
$\beta$ の値が物理の全セクターをパラメータ化する:

\begin{center}
\begin{tabular}{lll}
\toprule
$\beta$ & $\zeta(\beta)$ / $\zeta_{\neg 2}(\beta)$ & 物理 \\
\midrule
$\beta > 1$ & 収束級数 & 摂動的 QFT（Kreimer-Connes） \\
$\beta = 1$ & 極 & 相転移（Hagedorn） \\
$0 < \beta < 1$ & 対称性回復 & 大統一 \\
$\beta = 0$ & $\zeta(0) = -1/2$ & トポロジカル \\
$\beta = -1$ & $\zeta_{\neg 2}(-1) = +1/12$ & \textbf{ダークエネルギー (1D)} \\
$\beta = -3$ & $\zeta_{\neg 2}(-3) = -7/120$ & \textbf{負エネルギー (3D)} \\
\bottomrule
\end{tabular}
\end{center}

\noindent
関数等式 $\xi(\beta) = \xi(1-\beta)$ により:
\begin{itemize}
\item $\beta = -1 \leftrightarrow \beta = 2$:
  ダークエネルギー $\leftrightarrow$ $\zeta(2) = \pi^2/6$（摂動的QFT）
\item $\beta = -3 \leftrightarrow \beta = 4$:
  負エネルギー $\leftrightarrow$ $\zeta(4) = \pi^4/90$
\end{itemize}


%======================================================================
\section{CKN 不要: $\Omega_\Lambda = (d/\mathrm{cd})\,|\xi_{\neg 2}(-1)|$}
\label{sec:omega_formula}
%======================================================================

従来の $\Omega_\Lambda$ 導出では Cohen-Kaplan-Nelson ホログラフィック境界を
外部入力として必要としていた．BC 系の解析により，この仮定は不要になった:

\begin{tcolorbox}[colback=blue!5,colframe=blue!50!black,title=新公式]
\begin{equation}
  \Omega_\Lambda = \frac{d}{\mathrm{cd}} \times |\xi_{\neg 2}(-1)|
  = \frac{4}{3} \times \frac{\pi}{6} = \frac{2\pi}{9} \approx 0.698
\end{equation}
ここで:
\begin{itemize}
\item $d = 4$: 時空次元 ($= \mathrm{cd} + 1$, WDW 時間方向)
\item $\mathrm{cd} = 3$: $\mathrm{cd}(\mathrm{Spec}(\mathbb{Z}))$
  (Artin-Verdier 定理)
\item $\xi_{\neg 2}(-1) = -\pi/6$: completed ゼータ関数の $p=2$ 除去版
\end{itemize}
\end{tcolorbox}

\noindent
$(8\pi/3)$ という ``mysterious'' な前因子は:
\begin{equation}
  \frac{8\pi}{3} = \frac{d}{\mathrm{cd}} \times 2\pi
  = (\text{次元比}) \times (\text{アルキメデス}\,\Gamma\text{因子})
\end{equation}
と分解される．$2\pi$ は $\mathbb{Q}$ の無限素点（実数体 $\mathbb{R}$）の
寄与であり，$1/12$ は有限素点（素数 $p$）の寄与．
completed ゼータ $\xi$ はこの両方を統合した\textbf{自然な対象}である．

\subsection{三因子分解の統一}

\begin{center}
\begin{tabular}{lll}
\toprule
旧分解 & 新分解 & 起源 \\
\midrule
$4/3$ (起源不明) & $d/\mathrm{cd}$ & 時空の次元構造 \\
$2\pi$ (``ホログラフィック'') & $\xi$ の $\Gamma$ 因子 & $\mathbb{Q}$ の無限素点 \\
$1/12$ ($\zeta(-1)$) & $\zeta_{\neg 2}(-1)$ & $\mathbb{Q}$ の有限素点 ($p=2$ 除去) \\
\midrule
& $\rightarrow$ 2因子に統合 & $\Omega_\Lambda = (d/\mathrm{cd})\,|\xi_{\neg 2}(-1)|$ \\
\bottomrule
\end{tabular}
\end{center}


%======================================================================
\section{実験への影響}
%======================================================================

\subsection{何が変わったか}

\begin{tcolorbox}[colback=gray!10,colframe=gray!50!black,title=前回 (v1)]
\textbf{理解}: DN 境界は $\zeta_{\neg 2}(-1) > 0$ により斥力を生む．\\
\textbf{関心}: 符号反転の検出．\\
\textbf{問題}: ``負のエネルギー'' との接続が曖昧．
\end{tcolorbox}

\begin{tcolorbox}[colback=green!5,colframe=green!60!black,title=現在 (v2)]
\textbf{理解}: DN 境界の 3D cavity は $\zeta_{\neg 2}(-3) < 0$ により
\textbf{負のエネルギー密度}を生む．\\
\textbf{関心}: \textbf{3D 負エネルギーの直接検出}．\\
\textbf{接続}: これが Alcubierre ワープに必要な exotic matter の
算術的起源である．
\end{tcolorbox}

\subsection{1D vs 3D: 実験装置の意味が変わった}

\begin{center}
\begin{tabular}{lll}
\toprule
& 1D (平行板 Casimir) & 3D (BAW cavity) \\
\midrule
ゼータ値 & $\zeta_{\neg 2}(-1) = +1/12$ & $\zeta_{\neg 2}(-3) = -7/120$ \\
符号 & 正（斥力） & \textbf{負（負エネルギー密度）} \\
Casimir 力 & 斥力（Boyer 型） & 引力 \\
物理的意味 & ダークエネルギーの実験室版 & \textbf{exotic matter の実験室版} \\
ワープとの関連 & 推進（斥力） & \textbf{ワープ（負エネルギー）} \\
\bottomrule
\end{tabular}
\end{center}

\textbf{BAW 共振器は 3D cavity}であり，測定しているのは
$\zeta_{\neg 2}(-3)$ である．
これが負であれば，\textbf{$p=2$ ミュートによる
負のエネルギー密度生成が実験的に確認される}．

\subsection{なぜ BAW が 3D か}

BAW (Bulk Acoustic Wave) 共振器の音響モード:
\begin{equation}
  \omega_{n_x, n_y, n_z} = v_s \sqrt{(n_x \pi/L_x)^2 + (n_y \pi/L_y)^2 + (n_z \pi/L_z)^2}
\end{equation}

厚み方向 ($z$) に加え，横方向 ($x, y$) にもモードが立つ．
したがって零点エネルギーの正則化には 3D ゼータ値が関与する:
\begin{equation}
  E_{\mathrm{vac}}^{\mathrm{3D}} \propto \zeta_{\mathcal{H}}(-3)
\end{equation}

DN 境界 (上面 Dirichlet, 下面 Neumann) のとき:
\begin{equation}
  \zeta_{\mathcal{H}}(-3) \to \zeta_{\neg 2}(-3) = -\frac{7}{120} < 0
\end{equation}

\textbf{BAW 共振器は 3D 負エネルギーの自然な検出器である．}

%======================================================================
\section{更新された実験計画}
%======================================================================

\subsection{Phase 0: 発注 (¥130,000)}

変更なし．水晶 BAW + SMR/FBAR を発注．
\textbf{追加}: 実験の目的説明を「3D 負エネルギー検出」に更新．

\subsection{Phase 1: 基礎測定 + 3D 符号テスト (1ヶ月)}

\begin{enumerate}
\item \textbf{Week 1-2}: 組み立て + モードスペクトル確認
\item \textbf{Week 3}: \textbf{3D 符号テスト}（最重要）
  \begin{itemize}
  \item DN 共振器と DD 共振器の温度応答比較
  \item $\zeta_{\neg 2}(-3) < 0$ なら DN の温度応答は DD と\textbf{同じ方向}
    （どちらも負エネルギー寄与）
  \item ただし係数が異なる: DD は $\zeta(-3) = +1/120$（正），
    DN は $\zeta_{\neg 2}(-3) = -7/120$（負）
  \item 差分: $\Delta\zeta = -7/120 - 1/120 = -8/120 = -1/15$
  \end{itemize}
\item \textbf{Week 4}: データ解析
\end{enumerate}

\subsection{Phase 2: 88\,$\mu$m 焦点 + 次元テスト (2ヶ月)}

\begin{center}
\begin{tabular}{llll}
\toprule
厚さ & 周波数 & 次元 & テスト内容 \\
\midrule
\textbf{83\,$\mu$m} & \textbf{20\,MHz} & 3D & $\ell_\Lambda$ 近傍 (\textbf{最重要}) \\
330\,$\mu$m & 5\,MHz & 3D & プラトー内 \\
17\,$\mu$m & 100\,MHz & 準2D & 次元クロスオーバー \\
\bottomrule
\end{tabular}
\end{center}

\textbf{新テスト}: 極薄サンプル ($< 10\,\mu$m) で準 2D 化した場合，
$\zeta_{\neg 2}(-1) > 0$ に近づくはず → \textbf{符号が反転する厚さ}の特定．

\begin{tcolorbox}[colback=yellow!5,colframe=yellow!60!black,title=次元クロスオーバー予測]
\textbf{3D} (厚い BAW): $E_\mathrm{vac} \propto \zeta_{\neg 2}(-3) < 0$ （負エネルギー）\\
\textbf{準2D} (極薄 BAW): $E_\mathrm{vac} \propto \zeta_{\neg 2}(-1) > 0$ （正エネルギー）\\[4pt]
$\Rightarrow$ ある厚さで\textbf{符号がゼロを横切る}はず．
この ``critical thickness'' の測定は理論の強力なテスト．
\end{tcolorbox}

\subsection{Phase 3: 分岐 (3--6ヶ月)}

\begin{itemize}
\item \textbf{分岐 A (3D 負符号 confirm)}:
  \begin{itemize}
  \item 液体窒素冷却で精密測定
  \item 次元クロスオーバー厚の精密決定
  \item 係数 $7/120$ の精密値測定
  \item \textbf{``Negative vacuum energy in 3D DN cavity'' 論文投稿}
  \end{itemize}
\item \textbf{分岐 B (符号曖昧)}: 温度差分法の改良，系統誤差の削減
\item \textbf{分岐 C (null)}: 理論の棄却条件の検討
\end{itemize}

\subsection{Phase 4: スケールアップ (6ヶ月--1年)}

もし Phase 3A が成功した場合:
\begin{itemize}
\item \textbf{cavity 体積の増大}: 負エネルギー密度 $\times$ 体積 = 総負エネルギー
\item \textbf{cavity アレイ}: 多数の DN cavity の集積
\item \textbf{超伝導 cavity}: Q 値向上 → 負エネルギーのコヒーレンス向上
\item \textbf{共同研究提案}: ``Arithmetic vacuum engineering for exotic matter''
\end{itemize}


%======================================================================
\section{ワープドライブへの含意}
%======================================================================

\subsection{エネルギースケールの評価}

3D cavity でのDN Casimir エネルギー密度:
\begin{equation}
  \rho_{\mathrm{DN}}^{\mathrm{3D}} = \frac{\hbar\omega_0}{2V}\,\zeta_{\neg 2}(-3)
  = -\frac{7\hbar\omega_0}{240\,V}
\end{equation}

83\,$\mu$m BAW ($\omega_0 = 2\pi \times 20\,\mathrm{MHz}$, $V \sim (1\,\mathrm{mm})^2 \times 83\,\mu\mathrm{m}$):
\begin{equation}
  \rho_{\mathrm{DN}} \sim -10^{-18}\;\mathrm{J/m^3}
  \quad (\text{ダークエネルギー密度} \sim 10^{-10}\;\mathrm{J/m^3} \text{と比較})
\end{equation}

現時点では8桁不足．しかし:
\begin{itemize}
\item cavity サイズ縮小 ($d \to \mathrm{nm}$): $\rho \propto d^{-4}$ で増大
\item 超伝導 cavity の Q 値: $10^6$ 以上で coherent enhancement
\item アレイ化: 体積的スケーリング
\end{itemize}

\subsection{ロードマップ}

\begin{center}
\begin{tabular}{lll}
\toprule
段階 & 目標 & 必要技術 \\
\midrule
1. 原理実証 & $\zeta_{\neg 2}(-3) < 0$ の符号確認 & BAW ¥130k \\
2. 定量測定 & $7/120$ の精密値 & 低温 BAW \\
3. 次元制御 & 3D→2D クロスオーバー & 極薄 BAW \\
4. スケールアップ & ダークエネルギー密度に接近 & ナノ gap + cryo \\
5. 工学的応用 & マクロな負エネルギー領域 & 超伝導 cavity array \\
\bottomrule
\end{tabular}
\end{center}


%======================================================================
\section{成功判定基準}
%======================================================================

\begin{center}
\begin{tabular}{llp{7cm}}
\toprule
結果 & 判定 & 含意 \\
\midrule
3D DN で負符号 & \textbf{理論 confirm} &
  $p=2$ ミュートが 3D で負エネルギーを生むことの実証．
  ワープドライブの原理的可能性を支持 \\
3D/2D で符号反転 & \textbf{強い confirm} &
  次元依存性 $\zeta_{\neg 2}(-1) > 0$, $\zeta_{\neg 2}(-3) < 0$ の
  両方を確認 \\
88\,$\mu$m で異常 & \textbf{宇宙論的含意} &
  $\Omega_\Lambda = 2\pi/9$ との独立な接続 \\
全て null & \textbf{理論棄却} &
  Arithmetic landscape の実験的否定 \\
\bottomrule
\end{tabular}
\end{center}

%======================================================================
\section{Euclid/DESI との連携}
%======================================================================

\begin{center}
\begin{tabular}{lll}
\toprule
時期 & 実験室 (BAW) & 宇宙論 \\
\midrule
2026 Q2 & Phase 0 発注 & DESI Y1 データ \\
2026 Q3--Q4 & Phase 1--2: 3D 符号テスト & \\
2027 H1 & Phase 3: 次元クロスオーバー & Euclid 初期データ \\
2027 H2 & Phase 4: スケールアップ & DESI Y3 \\
2028 & 精密測定 & Euclid $\Omega_\Lambda$ 0.5\% 精度 \\
\bottomrule
\end{tabular}
\end{center}

\noindent
2028年に両 front の結果が揃う:
\begin{itemize}
\item 宇宙論: $\zeta_{\neg 2}(-1) = +1/12$ を $\Omega_\Lambda$ で検証
\item 実験室: $\zeta_{\neg 2}(-3) = -7/120$ を BAW cavity で検証
\item \textbf{同じ L 関数の 2 つの独立な evaluation}
\end{itemize}


%======================================================================
\section{論文パイプライン}
%======================================================================

\begin{enumerate}
\item \textbf{Paper 1}: ``Sign of vacuum energy in 3D DN acoustic cavity''
  → 符号検出で即投稿 (APL / Phys.\ Rev.\ B)

\item \textbf{Paper 2}: ``Dimensional crossover of vacuum energy sign
  in prime-muted cavities''
  → 3D/2D 符号反転 (PRL 候補)

\item \textbf{Paper 3}: ``Vacuum energy at the dark energy length scale
  $\ell_\Lambda = 88\,\mu$m''
  → 88\,$\mu$m 異常 (PRL 候補)

\item \textbf{Paper 4}: ``Bost-Connes negative temperature and the
  arithmetic origin of dark energy''
  → 理論論文 (Phys.\ Rev.\ D)

\item \textbf{Paper 5}: ``Laboratory evidence for the arithmetic vacuum
  $\zeta_{\neg 2}$''
  → 全データ統合 (Nature 候補)
\end{enumerate}


%======================================================================
\section{結論}
%======================================================================

\begin{tcolorbox}[colback=green!5,colframe=green!60!black]
\textbf{核心メッセージ}:

宇宙のダークエネルギー（$s=-1$, 正）と，
ワープに必要な negative energy（$s=-3$, 負）は，
\textbf{同じ素数ミュート $\zeta_{\neg 2}$ の異なる次元での評価}
に過ぎない．

¥130,000 の BAW 共振器で測定するのは $\zeta_{\neg 2}(-3) = -7/120$
の\textbf{符号}．
もしこれが負であれば:
\begin{enumerate}
\item ダークエネルギーの算術的起源が支持される
\item 同じ物理メカニズムで negative energy が生成可能であることが実証される
\item Alcubierre ワープの\textbf{原理的障壁}（exotic matter の存在）
  に対する実験的回答となる
\end{enumerate}

Euclid が $\Omega_\Lambda = 2\pi/9$ を宇宙論で検証し，
我々が $\zeta_{\neg 2}(-3) < 0$ を実験室で検証する．

2028年，同じ L 関数の 2 つの顔が揃う．
\end{tcolorbox}

\vspace{1em}
\hrule
\vspace{0.5em}
\small
\noindent
理論的基盤 (更新):
\begin{itemize}
\item \texttt{bc\_deep2.py}: $\Omega_\Lambda = (d/\mathrm{cd})\,|\xi_{\neg 2}(-1)|$ 導出
\item \texttt{bost\_connes\_negative.py}: フェルミオン/ボソン転換，Berry 位相
\item \texttt{bc\_pari\_deep.py}: PARI/GP による高精度検証
\item \texttt{bc\_sage\_minimal.sage}: SageMath によるモジュラー形式・K理論検証
\item \texttt{arithmetic\_vacuum\_engineering.tex}: 素数チャネルミュートの原論文
\end{itemize}

\end{document}
